\chapter{Experimentos e Resultados}
\label{cap:experimentos}
\label{cap:resultados}

\lc{Sem baseline: executar um analisador estático ou os testes existentes nestes módulos.}

Este capítulo apresenta a avaliação empírica da metodologia descrita no
Capítulo~\ref{cap:metodologia}, aplicada a 12 módulos do PX4 Autopilot
\artigo{(ramo \textit{main}, \textit{commit}~\texttt{8be2590c15})}.
\corr{As Seções~\ref{sec:config-experimental} a~\ref{sec:severidade} tratam
dessa primeira campanha; os resultados da segunda, com dez módulos
complementares e a validação em bancada, estão na
Seção~\ref{sec:resultados-segunda-campanha}.}

\revisao{O resultado principal da campanha é a evidência limitada em
nível de modelo: dez dos doze módulos, cobrindo as principais classes de
entrada não confiável, verificam-se livres de violações de segurança de memória
e de \textit{overflow}} \corr{das classes habilitadas, nos caminhos alcançados
sob $k = 10$\artigo{; a reexecução com as asserções de desenrolamento ligadas
mostrou que nenhum laço desses dez \textit{harnesses} excede o limite
(Seção~\ref{sec:config-experimental})} --- qualificação cujo alcance exato é discutido na
Seção~\ref{sec:ameacas-validade}.} \revisao{O resultado secundário é
a triagem disciplinada de três contraexemplos, que levou a três conclusões
distintas: um falso positivo no contador de FIFO do BMI088, refutado \artigo{na
conferência contra o código de produção}; uma observação de temperatura no BMI088 recusada pelos
mantenedores; e um \textit{underflow} de inteiro sem sinal no driver GPS uBlox}
\artigo{,
proposto upstream como endurecimento defensivo (PR~\#27576) e posteriormente
mostrado inalcançável no driver de produção}\revisao{. Essa distinção entre os três
desfechos é ela própria um resultado: contraexemplos de BMC expõem estados
alcançáveis no modelo que o código de produção pode nunca alcançar, de modo que
a triagem e a revisão pela comunidade são companhias necessárias às provas.}

\section{Configuração Experimental}
\label{sec:config-experimental}

A verificação foi executada em um servidor Intel Core i7-13700K com 64~GB de
RAM sob Ubuntu 24.04, sobre as fontes do PX4
\artigo{(ramo \textit{main}, \textit{commit}~\texttt{8be2590c15}, de 04 de
junho de 2026, 2.935 \textit{commits} após a v1.16.0 e posterior ao lançamento
da v1.17.0)}. Os tempos
reportados são a \textbf{mediana de três execuções}.

O artefato completo --- harnesses, logs de execução e instruções de reprodução
--- está arquivado no Zenodo sob o DOI~\artigo{\texttt{10.5281/zenodo.23282236}
(versão \texttt{v1.2-sbesc}; DOI de conceito \texttt{10.5281/zenodo.21401958})} e
disponível em \url{https://github.com/marcela2881/dissertacao-px4-esbmc}.

\artigo{Todos os doze módulos foram verificados com ESBMC v8.1.0 utilizando
Boolector~3.2.4~\cite{Brummayer2009} como solver; os 36 logs da campanha
(três execuções por módulo) registram \texttt{ESBMC version 8.1.0} e
\texttt{Solving with solver Boolector 3.2.4}. O Boolector é também a
configuração sob a qual o reporte upstream do GPS (Issue~\#26866) foi
registrado; em execuções exploratórias, o Z3~\cite{deMoura2008} não produziu
veredito no harness do GPS dentro do \textit{timeout}
(Seção~\ref{sec:ameacas-validade}).} A invocação padrão foi:

\begin{lstlisting}[
  caption={Comando padrão de execução do ESBMC.},
  label={lst:comando-esbmc},
  numbers=left,
  frame=single,
  language=bash]
esbmc modulo.cpp --unwind 10 --no-unwinding-assertions \
  --overflow-check --default-solver boolector --timeout 120s
\end{lstlisting}

Três módulos exigiram ajuste. O MAVLink utilizou \texttt{--unwind~300}, pois
seus laços de \textit{parsing} são profundos. O Navigator/Mission utilizou
\texttt{--unwind~15}. O VTOL utilizou \texttt{--no-library}, que desabilita a
biblioteca C abstrata do ESBMC: \texttt{sinf} e \texttt{sqrtf} passam a
retornar valores não determinísticos, uma sobreaproximação sob a qual nenhuma
violação das demais verificações é perdida; em contrapartida, as
\artigo{verificações internas da biblioteca deixam de ser feitas}, de modo que o
veredito cobre \artigo{correspondentemente} menos propriedades. A aritmética de ponto
flutuante usa a codificação SMT padrão do ESBMC
(\texttt{--floatbv})~\cite{Gadelha2020}.

Como \texttt{--no-unwinding-assertions} está ativo, laços que excedem $k = 10$
são truncados silenciosamente --- consequência discutida na
Seção~\ref{sec:ameacas-validade} \artigo{e verificada a seguir}.

%% LC: \textbf{[IMPORTANTE]} RECOMENDAÇÃO TRANSVERSAL Nº 3 --- A QUE MAIS FORTALECE O RESULTADO PRINCIPAL, E CUSTA CERCA DE TRÊS MINUTOS DE MÁQUINA.
%% LC: O resultado principal declarado é ``dez dos doze módulos verificam-se livres de violações até o limite de desenrolamento''. Mas a campanha roda com --no-unwinding-assertions. A combinação tem um efeito que a dissertação reconhece com todas as letras nas Ameaças à Validade: laços que excedem k são truncados SILENCIOSAMENTE, ``de modo que um veredito SUCCESSFUL pode ocultar falsos negativos''. Ou seja: o resultado principal é enunciado sob a única configuração que o impede de ser conclusivo, e o próprio texto o admite.
%% LC: Basta re-executar os DEZ módulos sem violação COM asserções de desenrolamento (removendo --no-unwinding-assertions). Só há dois desfechos, e ambos melhoram o trabalho:
%% LC: \hangindent=1em\hangafter=1 - nenhum módulo dispara a asserção -> está PROVADO que nenhum laço excede k, e o veredito passa a ser conclusivo dentro do recorte;
%% LC: \hangindent=1em\hangafter=1 - algum módulo dispara -> descobre-se onde o k é insuficiente e eleva-se o limite naquele módulo. Como a campanha inteira leva ~160 s, o custo é desprezível.
%% LC: Alternativa, se algum laço não fechar: --k-induction nos módulos afetados.

\artigo{\textbf{Reexecução com as asserções de desenrolamento ligadas.} Para
retirar do resultado principal a dependência do truncamento silencioso, os doze
\textit{harnesses} foram reexecutados em 12/10/2026 com as mesmas opções da
Listagem~\ref{lst:comando-esbmc} e os mesmos ajustes por módulo, mas
\textbf{sem} \texttt{--no-unwinding-assertions}, com o ESBMC~8.1.0 e o
Boolector~3.2.4 recompilados e validados contra os logs de julho. A
Tabela~\ref{tab:unwinding} resume os desfechos.}

\begin{table}[H]
\centering
\caption{\artigo{Reexecução com as asserções de desenrolamento ligadas}}
\label{tab:unwinding}
\footnotesize
\begin{tabular}{|p{6.3cm}|p{3.0cm}|p{4.2cm}|}
\hline
\textbf{Módulos} & \textbf{Veredito} & \textbf{Propriedade violada} \\
\hline
\artigo{Barômetro, Battery Monitor, EKF2 Covariance, Expo RC, MAVLink Parser
($k = 300$), Magnetômetro, Navigator/Mission ($k = 15$), PWM Output, RC Update
e VTOL} & \artigo{\textit{sem violação}} & \artigo{---} \\
\hline
\artigo{IMU BMI088} & \artigo{\textit{violação}} & \artigo{a mesma da campanha
(candidato do FIFO)} \\
\hline
\artigo{GPS uBlox} & \artigo{\textit{violação}} & \artigo{asserção de
desenrolamento do laço do \texttt{memcpy}} \\
\hline
\end{tabular}
\end{table}

\artigo{Nos dez módulos sem violação, nenhuma asserção de desenrolamento
falha: nenhum laço desses \textit{harnesses} excede o limite usado, e o veredito
deixa de depender do truncamento silencioso. No IMU, a primeira violação
continua sendo a do candidato do FIFO. No GPS, a primeira falha passa a ser a
asserção de desenrolamento do laço do \texttt{memcpy}, que copia mais de dez
bytes: a violação da Tabela~\ref{tab:resultados} continua válida no modelo,
porque o BMC é \textit{sound}, mas a análise do GPS com $k = 10$ não cobre esse
laço por inteiro. Como controle, o MAVLink com $k = 10$ e as asserções ligadas
falha apenas na asserção de desenrolamento do \texttt{memcpy}, e com $k = 300$
passa. Os logs estão em \texttt{verificacao\_unwinding\_2026-10/} no
repositório de artefatos.}

\lc{Laço do GPS excede k = 10; reexecutar com k cobrindo 210 bytes; reportar resultado.}

%% LC: ``nove dos onze logs'' --- são doze módulos. Onze logs de quê? Se o BMI088 e o GPS geram um log cada e os dez restantes também, seriam doze. Se algum módulo compartilha log, dizer qual. Como está, o leitor tropeça num número que não fecha com nenhum outro do capítulo.

\revisao{Quanto à proveniência das versões:} \artigo{os 36 logs de
\texttt{medianas\_logs/} (três execuções de cada um dos doze módulos)
registram explicitamente \texttt{ESBMC version 8.1.0}, conforme a nota de
rodapé do artigo (``read from all twelve run logs'').} \revisao{Os contraexemplos do BMI088 e do GPS foram
originalmente reportados upstream (Issues~\#26865 e~\#26866) sob o ESBMC
v7.10.0; os tempos aqui reportados foram remedidos sob a configuração da
campanha em v8.1.0.}

%% [BANCA - Prof. Barreto] "Tres valores diferentes de --unwind. Nao deveria justificar
%% melhor? A disparidade e muito grande, de 10, que e o padrao, para 300."
\subsection{\revisao{Escolha do Limite de Desenrolamento}}
\label{subsec:escolha-k}

\revisao{A campanha utilizou três valores distintos de $k$, e a disparidade entre eles
--- sobretudo entre o padrão $k = 10$ e o $k = 300$ do MAVLink --- merece justificativa
explícita, por ser o parâmetro que mais afeta o alcance das conclusões.}

\revisao{O limite de desenrolamento governa um compromisso direto. Aumentá-lo amplia o
conjunto de execuções cobertas e, portanto, a força do veredito; em contrapartida, cada
iteração adicional replica as instruções do laço na fórmula SMT, de modo que o custo cresce
rapidamente e pode tornar a verificação intratável dentro do \textit{timeout}. Não existe,
assim, um valor universalmente correto: o valor adequado é o menor que cobre o comportamento
de interesse do módulo.}

\revisao{O critério adotado foi derivar $k$ da estrutura do laço dominante de cada
módulo, e não arbitrá-lo:}

\begin{itemize}
\item \revisao{\textbf{$k = 10$ (padrão, nove módulos).} Nesses módulos, os laços
percorrem estruturas de dimensão pequena e fixa --- canais de rádio, saídas de PWM, elementos
de matriz de covariância --- todos com menos de dez iterações no pior caso. O valor cobre o
laço inteiro, e aumentá-lo não acrescentaria execuções novas.} \artigo{A reexecução com as
asserções de desenrolamento ligadas confirma isso: nenhuma delas falha nesses
\textit{harnesses} (Tabela~\ref{tab:unwinding}).}

\item \revisao{\textbf{$k = 15$ (Navigator/Mission).} O laço itera sobre
\textit{waypoints} da missão carregada\artigo{, e com $k = 15$ nenhuma asserção de
desenrolamento falha no \textit{harness} (Tabela~\ref{tab:unwinding}).}}

\item \revisao{\textbf{$k = 300$ (MAVLink Parser).} Aqui o valor não é arbitrário nem
uma tentativa de aumentar cobertura por precaução: o \texttt{memcpy} \artigo{do
\textit{harness}, que copia para o buffer de transmissão,} itera
proporcionalmente ao tamanho do pacote, e o buffer \artigo{é de 280 bytes, o tamanho
máximo de pacote}. Qualquer $k$
inferior a 280 deixaria de cobrir pacotes de tamanho máximo\artigo{. A reexecução com
as asserções ligadas confirma isso: com $k = 10$, a única falha é a asserção de
desenrolamento do \texttt{memcpy} --- um resultado inconclusivo por profundidade
insuficiente, e não um defeito ---, e com $k = 300$ o \textit{harness} passa
(Seção~\ref{sec:config-experimental}).}
O valor 300 é o menor múltiplo redondo acima do
tamanho do buffer.}
\end{itemize}

\revisao{Segue-se disso que a disparidade entre 10 e 300 não reflete inconsistência de
critério, mas heterogeneidade dos módulos: o MAVLink é o único cujo laço dominante percorre um
buffer de centenas de bytes. Duas consequências devem ser registradas. Primeiro, os totais de
VCC por módulo misturam limites diferentes e, portanto, não são diretamente comparáveis entre
si. Segundo, o procedimento é reativo por construção --- parte-se do padrão e eleva-se $k$
quando o verificador sinaliza insuficiência --- de modo que um módulo cujo laço exceda o limite
sem disparar asserção de desenrolamento poderia passar despercebido, limitação agravada pelo
uso de \texttt{--no-unwinding-assertions} e discutida na
Seção~\ref{sec:ameacas-validade}.} \artigo{A reexecução com as asserções ligadas
(Tabela~\ref{tab:unwinding}) afasta esse risco nos dez módulos sem violação.}

\section{Questões Experimentais}
\label{sec:objetivos-verificacao}

A avaliação responde a três questões experimentais (QE). \corr{Para cada uma
indicam-se a métrica observada e o critério adotado para considerá-la
respondida.}

\begin{itemize}
    \item \textbf{QE1:} Qual é a latência de verificação por módulo em hardware
    de prateleira? \corr{\textit{Métrica:} tempo do procedimento de decisão,
    mediana de três execuções, com mínimo e máximo.
    \textit{Critério:} todo módulo conclui dentro do \textit{timeout} de
    $120$\,s.}
    \item \textbf{QE2:} O BMC expõe violações alcançáveis no modelo em código
    que os testes unitários existentes do PX4 não referenciam?
    \corr{\textit{Métrica:} número de contraexemplos, desfecho de cada um após
    a triagem e existência de teste unitário que referencie a função afetada.
    \textit{Critério:} ao menos um contraexemplo que sobreviva à triagem como
    violação alcançável no modelo, em função sem teste unitário direto.}
    \item \textbf{QE3:} Quais configurações por módulo são necessárias para
    tornar cada módulo tratável? \corr{\textit{Métrica:} número de módulos que
    exigem opções diferentes da invocação padrão, e quais opções.
    \textit{Critério:} questão descritiva, sem limiar de aceitação.}
\end{itemize}

\corr{Os critérios foram fixados depois da campanha, ao redigir este capítulo;
servem para tornar explícito o que cada resposta afirma, e não como hipóteses
registradas previamente.}

\section{Módulos Verificados}
\label{sec:modulos-verificados}

Os módulos foram priorizados pelo impacto na segurança ao longo do caminho de
sensoriamento até atuação e pela viabilidade de isolamento por \textit{stubs}.
Foram verificados 12 módulos por meio de \textit{harnesses} construídos com assistência de IA generativa e revisados pela autora,
totalizando aproximadamente $2.800$ linhas, medidas com \texttt{wc -l} sobre os
arquivos \texttt{test\_*.cpp} do repositório de artefatos.

\revisao{Em vez de reportar percentual de cobertura de linhas --- que
seria enganoso, já que cada harness abstrai a maior parte do seu módulo e
explora apenas caminhos alcançáveis sob $k = 10$ --- a cobertura é
caracterizada por \textbf{classe de entrada}, conforme a última coluna da
Tabela~\ref{tab:modulos}.}

\begin{table}[H]
\centering
\caption[Módulos verificados na primeira campanha]{Módulos verificados: tamanho do alvo, tamanho do harness, localização
no código-fonte do PX4 e classe de entrada exercitada}
\label{tab:modulos}
\footnotesize
\begin{tabular}{|l|l|r|r|l|}
\hline
\textbf{Módulo} & \textbf{Caminho (\texttt{src/})} & \textbf{Alvo} &
\textbf{Harness} & \textbf{Classe de entrada} \\
\hline
GPS uBlox         & \texttt{drivers/gps}          & \artigo{1712} & 294 & Sensor serial \\
IMU BMI088        & \texttt{drivers/imu/.../bmi088} & \artigo{700} & 319 & Sensor serial \\
MAVLink           & \texttt{modules/mavlink}      & \artigo{3624} &  82 & Telemetria \\
Expo RC           & \texttt{lib/mathlib/math}     &  \artigo{303} & 129 & Rádio \\
Battery Monitor   & \texttt{lib/battery}          &  \artigo{449} & 242 & Sensor de energia \\
PWM Output        & \texttt{drivers/pwm\_out}     &  337 & 208 & Atuação \\
RC Update         & \texttt{lib/rc}               & \artigo{769} & 253 & Rádio \\
Navigator         & \texttt{modules/navigator}    & \artigo{497} & 291 & Lógica de missão \\
EKF2 Covariance   & \texttt{modules/ekf2/EKF}     &  385 & 296 & Estimação \\
Magnetômetro      & \texttt{modules/sensors}      &  792 & 239 & \corr{Agregação de sensores} \\
Barômetro         & \texttt{modules/sensors}      &  \artigo{246} & 278 & \corr{Agregação de sensores} \\
VTOL              & \texttt{modules/vtol\_att\_control} & \artigo{591} & 146 & Atuação \\
\hline
\end{tabular}

\footnotesize
Alvo = LoC \artigo{do(s) arquivo(s)-fonte do PX4 que o harness extrai (\textit{commit} fixado)}; Harness = LoC do harness (ambos por \texttt{wc -l}).
\end{table}

\lc{Acrescentar coluna: linhas do PX4 efetivamente extraídas por harness, separadas do código de modelo.}
\corr{Os arquivos-alvo somam $10.405$ linhas, cerca de $1{,}5\,\%$ das $710.259$
linhas de C/C++ em \texttt{src/} no \textit{commit} analisado ($2.710$ arquivos,
contados com \texttt{wc -l}). Os alvos do Magnetômetro e do Barômetro estão em
\texttt{modules/sensors}: são módulos que agregam e selecionam leituras já
publicadas pelos drivers, e não drivers de barramento serial, razão pela qual
aparecem como ``agregação de sensores''.}

\section{Resultados Gerais}
\label{sec:resultados-gerais}

A Tabela~\ref{tab:resultados} sumariza os desfechos. Dez dos doze módulos
verificam-se sem violação até o limite\artigo{; os dois restantes (IMU BMI088 e
GPS uBlox) param na primeira violação encontrada. Como no artigo derivado, as
contagens de condições de verificação (VCCs) por \textit{harness} não são
tabeladas: elas misturam limites de desenrolamento diferentes (10, 15 e 300) e,
nos dois módulos com violação, cobrem só o caminho até a primeira violação. Os
valores constam dos logs do repositório de artefatos.}

\begin{table}[H]
\centering
\caption{Resultados da verificação para os 12 módulos do PX4}
\label{tab:resultados}
\footnotesize
\setlength{\tabcolsep}{4pt}
\begin{tabular}{|l|l|c|r|l|r|r|}
\hline
\textbf{Módulo} & \textbf{Desfecho} &
\textbf{VAM}$^\P$ & \textbf{Unw.} & \textbf{Solver} & \textbf{Tempo}$^\|$ & \corr{\textbf{Decisão}$^{\star}$ (mín--máx)} \\
\hline
IMU BMI088        & \corr{\textit{contraex.\ (triado)}} & 0$^\ddagger$ &  10 & \artigo{Bool.} & \artigo{0,1\,s}$^{\#}$ & \corr{0,122 (0,121--0,122)} \\
GPS uBlox         & \textit{violação}     & 1 &  10 & \artigo{Bool.} & 64,9\,s$^{\#}$ & \corr{64,857 (64,765--64,874)} \\
MAVLink Parser    & \textit{sem violação} & 0 & 300 & \artigo{Bool.} & 49,2\,s & \corr{46,970 (46,916--47,083)} \\
Expo RC           & \textit{sem violação} & 0 &  10 & \artigo{Bool.} &  4,9\,s & \corr{4,852 (4,840--4,866)} \\
Battery Monitor   & \textit{sem violação} & 0 &  10 & \artigo{Bool.} &  6,7\,s & \corr{6,727 (6,703--6,731)} \\
PWM Output        & \textit{sem violação} & 0 &  10 & \artigo{Bool.} &  0,2\,s & \corr{0,156 (0,156--0,156)} \\
RC Update         & \textit{sem violação} & 0 &  10 & \artigo{Bool.} &  9,5\,s & \corr{9,457 (9,384--9,534)} \\
Navigator/Mission & \textit{sem violação} & 0 &  15 & \artigo{Bool.} &  0,1\,s & \corr{0,095 (0,095--0,095)} \\
EKF2 Covariance   & \textit{sem violação} & 0 &  10 & \artigo{Bool.} & 12,2\,s & \corr{12,200 (12,181--12,228)} \\
Magnetômetro      & \textit{sem violação} & 0 &  10 & \artigo{Bool.} &  6,6\,s & \corr{6,621 (6,616--6,629)} \\
Barômetro         & \textit{sem violação} & 0 &  10 & \artigo{Bool.} &  0,4\,s & \corr{0,360 (0,351--0,363)} \\
VTOL$^\S$         & \textit{sem violação} & 0 &  10 & \artigo{Bool.} &  4,0\,s & \corr{3,942 (3,911--3,943)} \\
\hline
\textbf{Total} & \corr{\textbf{10 sem violação, 1 violação,}} & \textbf{1} &
--- & --- & \artigo{\textbf{$\approx$94\,s}$^{\dagger}$} & \corr{\textbf{156,359}} \\
               & \corr{\textbf{1 contraex.\ refutado/recusado}} & & & & & \\
\hline
\end{tabular}

\footnotesize
\begin{flushleft}
$^\P$VAM = violações alcançáveis no modelo (\textbf{não} defeitos
confirmados).\\
\artigo{Bool.\ = Boolector~3.2.4, com ESBMC v8.1.0, em todos os módulos.}\\
$^\ddagger$\corr{Dois contraexemplos, ambos triados como não-defeito:
FIFO, falso positivo; temperatura, recusada upstream.} \artigo{O tempo da
tabela é o do candidato FIFO (\texttt{test\_choice = 2}), em que as três
execuções param (0,122; 0,122; 0,121\,s); o contraexemplo de temperatura vem de
uma execução separada (0,272\,s), Seção~\ref{sec:caso-imu}.}\\
$^{\dagger}$\artigo{93,8\,s nos módulos que concluem (dez sem violação e o
IMU); 158,7\,s incluindo o GPS (64,9\,s).}\\
$^\S$VTOL com \texttt{--no-library}: \artigo{verificações internas da
biblioteca não são feitas.}\\
$^\|$Mediana de três execuções.\\
\corr{$^{\star}$Tempo do procedimento de decisão (\textit{Runtime decision procedure}), em segundos: mediana, mínimo e máximo das três execuções. É a única medida registrada nos logs dos doze módulos.}\\
$^{\#}$\artigo{Módulos com violação param na primeira violação; tempo do
procedimento de decisão (\textit{Runtime decision procedure}); nos demais,
\textit{BMC program time}.}
\end{flushleft}
\end{table}

O tempo máximo foi de $64,9$\,s (GPS) e o mínimo de $0,1$\,s
\artigo{(Navigator/Mission e IMU BMI088)}.

\corr{A coluna \textit{Tempo} reproduz a Tabela~II do artigo e mistura duas
medidas, conforme a nota~$^{\#}$: os logs dos dois módulos com violação não
registram o tempo total do programa. Para permitir comparação uniforme, a
última coluna traz o tempo do procedimento de decisão dos doze módulos, com
mínimo e máximo; a dispersão entre execuções é inferior a $0{,}2$\,s em todos
os casos. O tempo de relógio de ponta a ponta não foi registrado nos logs da
campanha. A Figura~\ref{fig:visao-resultados} resume desfechos e tempos.}

\begin{figure}[!htb]
\centering
\begin{tikzpicture}
  \draw[palEdge] (0,0.3) -- (0,-4.97);
  \draw[palEdge, densely dotted] (0.000,0.3) -- (0.000,-4.97);
  \node[font=\scriptsize, anchor=north] at (0.000,-4.97) {0,01};
  \draw[palEdge, densely dotted] (2.250,0.3) -- (2.250,-4.97);
  \node[font=\scriptsize, anchor=north] at (2.250,-4.97) {0,1};
  \draw[palEdge, densely dotted] (4.500,0.3) -- (4.500,-4.97);
  \node[font=\scriptsize, anchor=north] at (4.500,-4.97) {1};
  \draw[palEdge, densely dotted] (6.750,0.3) -- (6.750,-4.97);
  \node[font=\scriptsize, anchor=north] at (6.750,-4.97) {10};
  \draw[palEdge, densely dotted] (9.000,0.3) -- (9.000,-4.97);
  \node[font=\scriptsize, anchor=north] at (9.000,-4.97) {100};
  \node[anchor=east, font=\scriptsize] at (-0.1,-0.00) {GPS uBlox};
  \draw[fill=palStrong, draw=palEdgeDark, thick] (0,-0.14) rectangle (8.577,0.14);
  \node[anchor=west, font=\scriptsize] at (8.627,-0.00) {64,857\,s};
  \node[anchor=east, font=\scriptsize] at (-0.1,-0.42) {MAVLink Parser};
  \draw[fill=palFill, draw=palEdge] (0,-0.56) rectangle (8.262,-0.28);
  \node[anchor=west, font=\scriptsize] at (8.312,-0.42) {46,970\,s};
  \node[anchor=east, font=\scriptsize] at (-0.1,-0.84) {EKF2 Covariance};
  \draw[fill=palFill, draw=palEdge] (0,-0.98) rectangle (6.944,-0.70);
  \node[anchor=west, font=\scriptsize] at (6.994,-0.84) {12,200\,s};
  \node[anchor=east, font=\scriptsize] at (-0.1,-1.26) {RC Update};
  \draw[fill=palFill, draw=palEdge] (0,-1.40) rectangle (6.695,-1.12);
  \node[anchor=west, font=\scriptsize] at (6.745,-1.26) {9,457\,s};
  \node[anchor=east, font=\scriptsize] at (-0.1,-1.68) {Battery Monitor};
  \draw[fill=palFill, draw=palEdge] (0,-1.82) rectangle (6.363,-1.54);
  \node[anchor=west, font=\scriptsize] at (6.413,-1.68) {6,727\,s};
  \node[anchor=east, font=\scriptsize] at (-0.1,-2.10) {Magnetômetro};
  \draw[fill=palFill, draw=palEdge] (0,-2.24) rectangle (6.347,-1.96);
  \node[anchor=west, font=\scriptsize] at (6.397,-2.10) {6,621\,s};
  \node[anchor=east, font=\scriptsize] at (-0.1,-2.52) {Expo RC};
  \draw[fill=palFill, draw=palEdge] (0,-2.66) rectangle (6.043,-2.38);
  \node[anchor=west, font=\scriptsize] at (6.093,-2.52) {4,852\,s};
  \node[anchor=east, font=\scriptsize] at (-0.1,-2.94) {VTOL};
  \draw[fill=palFill, draw=palEdge] (0,-3.08) rectangle (5.840,-2.80);
  \node[anchor=west, font=\scriptsize] at (5.890,-2.94) {3,942\,s};
  \node[anchor=east, font=\scriptsize] at (-0.1,-3.36) {Barômetro};
  \draw[fill=palFill, draw=palEdge] (0,-3.50) rectangle (3.502,-3.22);
  \node[anchor=west, font=\scriptsize] at (3.552,-3.36) {0,360\,s};
  \node[anchor=east, font=\scriptsize] at (-0.1,-3.78) {PWM Output};
  \draw[fill=palFill, draw=palEdge] (0,-3.92) rectangle (2.685,-3.64);
  \node[anchor=west, font=\scriptsize] at (2.735,-3.78) {0,156\,s};
  \node[anchor=east, font=\scriptsize] at (-0.1,-4.20) {IMU BMI088};
  \draw[fill=palPaper, draw=palEdgeDark, dashed] (0,-4.34) rectangle (2.444,-4.06);
  \node[anchor=west, font=\scriptsize] at (2.494,-4.20) {0,122\,s};
  \node[anchor=east, font=\scriptsize] at (-0.1,-4.62) {Navigator/Mission};
  \draw[fill=palFill, draw=palEdge] (0,-4.76) rectangle (2.200,-4.48);
  \node[anchor=west, font=\scriptsize] at (2.250,-4.62) {0,095\,s};
  \draw[fill=palFill, draw=palEdge] (0,-6.37) rectangle (0.35,-6.17); \node[anchor=west, font=\scriptsize] at (0.4,-6.27) {sem violação};
  \draw[fill=palStrong, draw=palEdgeDark, thick] (2.6,-6.37) rectangle (2.95,-6.17); \node[anchor=west, font=\scriptsize] at (3.0,-6.27) {violação alcançável no modelo};
  \draw[fill=palPaper, draw=palEdgeDark, dashed] (7.2,-6.37) rectangle (7.55,-6.17); \node[anchor=west, font=\scriptsize] at (7.6,-6.27) {contraexemplo triado};
  \node[font=\scriptsize, anchor=north] at (4.50,-5.27) {tempo do procedimento de decisão (s, escala logarítmica)};
\end{tikzpicture}
\caption[Desfecho e tempo de decisão por módulo (primeira campanha)]{\corr{Desfecho e tempo do procedimento de decisão por módulo na primeira campanha (mediana de três execuções). Elaborado pela autora.}}
\label{fig:visao-resultados}
\end{figure}

\lc{Um main multiplexa vários testes; a primeira falha mascara as demais; executar cada teste separadamente.}

\begin{tcolorbox}[colback=blue!5!white, colframe=blue!60!black,
title=\textbf{Resposta à QE1 --- Latência de verificação}]
Todo módulo cronometrado verifica-se em menos de noventa segundos em hardware
de prateleira. Como cada um é verificado isoladamente por meio de um harness
pequeno, isso demonstra \textbf{viabilidade por módulo}, e não escalabilidade
de sistema completo.
\end{tcolorbox}

\revisao{Os módulos verificados não constituem uma amostra uniforme:
situam-se no caminho crítico de sensoriamento até atuação, onde bytes não
confiáveis de I2C, SPI e UART entram no \textit{flight stack} pela primeira
vez, cobrindo as principais classes de entrada não confiável --- barramentos
seriais de sensores, enlace de rádio e telemetria, atuação, estimação e lógica
de missão.} \corr{O que ficou de fora, porém, não é código periférico. No
\textit{commit} analisado, as $710.259$ linhas de C/C++ de \texttt{src/}
distribuem-se entre \texttt{modules} ($45{,}8\,\%$), \texttt{drivers}
($39{,}4\,\%$), \texttt{lib} ($11{,}4\,\%$) e \texttt{systemcmds}
($2{,}9\,\%$), e os maiores blocos não verificados incluem código crítico para
o voo: \texttt{drivers/uavcan} ($58.267$ linhas), \texttt{drivers/imu}
($46.127$), \texttt{modules/ekf2} ($33.516$, do qual só a rotina de covariância
foi verificada), \texttt{modules/commander} ($25.377$) e
\texttt{modules/navigator} ($17.613$). A campanha cobre, portanto, uma fração
pequena e escolhida do \textit{firmware}.} \revisao{As exceções críticas para voo (\textit{commander},
\textit{failsafe} e lógica de \textit{geofence}) são trabalho futuro
prioritário, e os barramentos CAN/UAVCAN, a ponte RTPS/uXRCE-DDS e os drivers
de fluxo óptico e telêmetro permanecem não verificados.}

\section{Validação Cruzada com o CBMC}
\label{sec:cbmc}

\revisao{Para avaliar a dependência dos resultados em relação à
ferramenta, o CBMC~5.95.1~\cite{Kroening2014} foi executado sobre os}
\corr{harnesses adaptados aos intrínsecos do CBMC
(\texttt{\_\_CPROVER\_assume} e \texttt{\_\_CPROVER\_assert} no lugar de
\texttt{\_\_ESBMC\_assume} e \texttt{\_\_ESBMC\_assert}), preservando
pré-condições e asserções,} \revisao{com opções equivalentes
(\texttt{--unwind 10}, verificações de \textit{overflow} e de limites, sem
asserções de desenrolamento) e o mesmo \textit{timeout} de $120$\,s.}
\corr{A adaptação é parte do experimento e está registrada no repositório de
artefatos}\artigo{, que contém o \textit{harness} do CBMC para o GPS
(\texttt{cbmc/}) e os logs do CBMC para o IMU e o GPS.}

\revisao{Para o IMU BMI088, as duas ferramentas concordaram. Para o GPS
uBlox, sob a pré-condição \texttt{len}~$\geq$~$190$
(\texttt{GPS\_DUMP\_DATA\_SIZE}~$-$~$10$, semeando o campo de comprimento
próximo à fronteira do buffer), ambas sinalizam a violação}
\artigo{, mas o CBMC apenas com limites reduzidos (\texttt{--unwind} 2, 3 e 5:
\textit{memcpy destination region writeable}); com \texttt{--unwind 10} ele não
terminou dentro de um limite de 30 minutos. O harness original do CBMC, com a
pré-condição \texttt{len}~$<$~$200$, retornou \textsc{verification successful},
porque a pré-condição exclui os estados violadores. O harness do ESBMC usa a
janela $\texttt{GPS\_DUMP\_DATA\_SIZE} - 5$. As duas ferramentas reportam a
mesma propriedade --- acesso fora dos limites dentro do \texttt{memcpy}
(\textit{dereference failure} no ESBMC); a cadeia CWE-191 $\rightarrow$
CWE-119/787 é uma classificação deste trabalho, e não saída das ferramentas.}

\corr{No IMU, o CBMC~5.95.1 reporta 2 de 6 propriedades violadas: as duas
asserções do \textit{harness} (faixa de temperatura e
\texttt{count <= FIFO\_SIZE}) falham, e as quatro verificações de
\textit{overflow} aritmético passam; o tempo do procedimento de decisão é
inferior a $0{,}01$\,s. São as mesmas duas propriedades violadas no ESBMC. O log
(\texttt{cbmc/resultado\_cbmc\_imu.txt}) não registra a linha de comando usada.
No GPS, o comando registrado é o da Listagem~\ref{lst:comando-cbmc}:}

\begin{lstlisting}[
  caption={\corr{Comando do CBMC para o \textit{harness} do GPS ($u \in \{2, 3, 5, 10\}$).}},
  label={lst:comando-cbmc},
  numbers=left,
  frame=single,
  language=bash]
cbmc test_gps_cbmc_fair.cpp --unwind $u --no-unwinding-assertions \
  --bounds-check --pointer-check \
  --signed-overflow-check --unsigned-overflow-check
\end{lstlisting}

\corr{As janelas de entrada dos dois \textit{harnesses} do GPS não são
idênticas: o do CBMC assume \texttt{len}~$\geq 190$ e o do ESBMC,
\texttt{len}~$\geq 195$. Os dois foram escritos em momentos diferentes e a
diferença não foi uniformizada. Ela não altera a comparação, porque os estados
que violam a propriedade (\texttt{len}~$> 200$) pertencem às duas janelas; uma
reexecução com a mesma janela nas duas ferramentas tornaria a comparação mais
limpa e fica registrada como pendência. As versões de compilador registradas
nos artefatos são as do \textit{replay} da segunda campanha (g++~13.3.0); o
ESBMC e o CBMC usam seus próprios \textit{front-ends}.}

\lc{Validação cruzada cobre dois módulos; executar CBMC também nos dez harnesses sem violação.}

\revisao{Cabe uma ressalva importante: isso \textbf{não} constitui
confirmação independente de alcançabilidade. Com a mesma pré-condição no
harness, a concordância corrobora apenas a codificação do problema, não que o
estado surja no driver de produção.} \corr{A ressalva é ainda mais forte quando
se considera que as duas execuções partem de codificações distintas do mesmo
problema: a concordância entre elas atesta que o defeito sobrevive à tradução
entre os dois conjuntos de intrínsecos, e nada além disso.}

\section{Estudo de Caso: Driver IMU BMI088}
\label{sec:caso-imu}

\corr{O driver do BMI088 (\artigo{$700$}~LoC no arquivo verificado, conforme a
Tabela~\ref{tab:modulos})} produziu dois candidatos.
\revisao{A triagem levou a conclusões diferentes para cada um --- e
nenhuma delas é a de defeito confirmado.}

\subsection{Temperatura fora da faixa operacional (contestado, recusado upstream)}
\label{subsec:imu-temperatura}

O driver não valida a temperatura antes da aritmética de compensação térmica; o
datasheet do BMI088~\cite{Bosch2021} restringe a operação à faixa de
$[-40,\,+85]\,^\circ$C (Equação~\ref{eq:prop_temp}). O ESBMC produziu
\texttt{temp\_msb = 63} e \texttt{temp\_lsb = 95}
(\texttt{Temp\_uint11 = 506}), resultando em $86{,}25\,^\circ$C.
\corr{A conversão especificada pelo fabricante é
$T = \texttt{Temp\_int11} \times 0{,}125 + 23$, de modo que
$506 \times 0{,}125 + 23 = 86{,}25\,^\circ$C; o traço literal consta do
Apêndice~\ref{apend:contraexemplos}.}

\corr{Se a observação fosse aceita como defeito, a classe em que se enquadraria
é a CWE-1284 (\textit{Improper Validation of Specified Quantity in Input}),
subclasse de CWE-20 (\textit{Improper Input Validation}): o driver consome uma
leitura fora da faixa do \textit{datasheet} sem verificação. Como a observação
foi recusada, o rótulo indica apenas a classe do cenário investigado, e não uma
vulnerabilidade confirmada.}

\revisao{Se isso constitui defeito é ponto contestado. Os mantenedores
observam que a temperatura não alimenta o caminho inercial, e que descartar
quadros da IMU diante de uma excursão transitória arrisca perda de controle em
plataformas com uma única IMU; tais exceções seriam melhor tratadas pelas
camadas de estimação a jusante, como a fusão de sensores do EKF2, do que por
paradas em nível de driver. Trata-se de um compromisso entre rigor no driver e
tolerância a falhas em nível de sistema.}

\revisao{Um \textit{patch} upstream que limitava a temperatura fora de
faixa (PX4~PR~\#27575) foi \textbf{rejeitado}: os mantenedores objetaram à
abordagem de limitar e avisar, e fecharam tanto o PR quanto a
Issue~\#26865\footnote{\url{https://github.com/PX4/PX4-Autopilot/issues/26865}}.
A observação é mantida nesta dissertação apenas como nota de projeto, \textbf{sem
consequência de segurança confirmada e sem remediação aceita}, ilustrando a
etapa de triagem da metodologia em vez de um defeito.} \corr{A
Figura~\ref{fig:ce-imu-temp} reproduz a saída do verificador.}

\begin{figure}[H]
    \centering
    \definecolor{ceFundo}{HTML}{F7F8FA}\definecolor{ceCinza}{HTML}{5A6070}%
    \definecolor{ceAmarelo}{HTML}{0B5CAD}\definecolor{ceVermelho}{HTML}{B3261E}%
    \begin{tcolorbox}[colback=ceFundo, colframe=black!45, boxrule=0.6pt, arc=2pt,
      left=8pt, right=8pt, top=6pt, bottom=6pt, width=0.95\linewidth,
      fontupper=\ttfamily\scriptsize\linespread{1.45}\selectfont]
    \textcolor{ceVermelho}{\textbf{[Counterexample]}}\\
    \textcolor{ceCinza}{State 2 -{}- test\_bmi088\_imu.cpp:120 -{}- test\_temperature\_calculation}\\
    \hspace*{1em}\fcolorbox{ceAmarelo}{blue!6}{\textcolor{ceAmarelo}{\textbf{temp\_msb = 63 (00111111)}}}\\
    \textcolor{ceCinza}{State 3 -{}- test\_bmi088\_imu.cpp:121 -{}- test\_temperature\_calculation}\\
    \hspace*{1em}\fcolorbox{ceAmarelo}{blue!6}{\textcolor{ceAmarelo}{\textbf{temp\_lsb = 95 (01011111)}}}\\
    \textcolor{ceCinza}{State 7 -{}- test\_bmi088\_imu.cpp:42 -{}- updateTemperature}\\
    \hspace*{1em}\fcolorbox{ceAmarelo}{blue!6}{\textcolor{ceAmarelo}{\textbf{Temp\_uint11 = 506 (00000001 11111010)}}}\\
    \textcolor{ceCinza}{State 11 -{}- test\_bmi088\_imu.cpp:51 -{}- updateTemperature}\\
    \hspace*{1em}\fcolorbox{ceAmarelo}{blue!6}{\textcolor{ceAmarelo}{\textbf{temperature = 8.625000e+1f (...00000000000)}}}\\
    \textcolor{ceVermelho}{\textbf{Violated property:}}\\
    \textcolor{ceVermelho}{\hspace*{1em}file test\_bmi088\_imu.cpp line 126 function test\_temperature\_calculation}\\
    \textcolor{ceVermelho}{\hspace*{1em}assertion temperature >= -40.0f \&\& temperature <= 85.0f}\\
    \textcolor{ceVermelho}{\hspace*{1em}temperature >= -40.000000f \&\& temperature <= 85.000000f}\\
    \textcolor{ceVermelho}{\textbf{VERIFICATION FAILED}}\\[-2pt]
    \textcolor{ceCinza}{\rule{\linewidth}{0.3pt}}\\
    \textcolor{ceAmarelo}{8.625000e+1f = 86,25\,$^{\circ}$C > 85\,$^{\circ}$C (limite operacional do datasheet BMI088)}\\
    \textcolor{ceCinza}{ESBMC 8.1.0 + Boolector 3.2.4 -{}- log: resultado\_temperatura\_8.1.log}\\
    \textcolor{ceCinza}{PX4 main, commit 8be2590c15 -{}- Issue \#26865 (fechada) / PR \#27575 (fechado)}
    \end{tcolorbox}
    \caption[Contraexemplo de temperatura do BMI088 (saída do ESBMC)]{\corr{Saída do ESBMC para o contraexemplo de temperatura do BMI088, montada a partir do log (ESBMC~v8.1.0, Boolector~3.2.4).}}
    \label{fig:ce-imu-temp}
\end{figure}

%% LC: \textbf{[EM ABERTO]} Conferir dois dados internos desta figura contra o texto: ela cita o datasheet como ``BST-BMI088-DS001 §5.3.7'', ao passo que o .bib registra ``BST-BMI088-DS000-19, rev. 1.9'' e a Tab. \ref{tab:proveniencia} repete essa segunda forma. São dois números de documento diferentes para a mesma fonte. Verificar qual é o correto no PDF da Bosch e uniformizar figura, .bib e tabela.
%% LC: \textbf{[AJUSTE]} Renomear figuras/imagematual.jpeg (usada na Fig. \ref{fig:ce-gps}). Sugestão: ce\_gps\_underflow.jpeg. Aproveitar para remover de figuras/ os arquivos que o LEIA-ME já identificou como contraditórios com o texto atual: bug1\_temperatura.jpeg (mostra 512 -> 87,0 $^\circ$C), bug\_gps.png (input\_len=1, len=216) e bug2\_fifo.jpeg (count=1032). Enquanto estiverem no .zip entregue, são um risco: basta alguém abrir a pasta.

\subsection{Contador do FIFO (falso positivo)}
\label{subsec:imu-fifo}

\revisao{O ESBMC produziu} \artigo{\texttt{count = 9088}} \revisao{bytes, mais de oito
vezes o FIFO de $1.024$ bytes, o que sugeria CWE-190 e CWE-787. A}
\artigo{conferência} \revisao{desse contraexemplo} \artigo{contra o código} \revisao{do driver de produção \textbf{não
instrumentado} revelou um \textbf{falso positivo}.}

\revisao{O driver liberado protege a leitura em três níveis: o contador
de bytes é comparado ao tamanho do FIFO e reiniciado sempre que
\texttt{fifo\_byte\_counter}~$\geq$~\texttt{FIFO::SIZE}; a contagem de amostras
é limitada por \texttt{FIFO\_MAX\_SAMPLES}; e o comprimento da transferência é
restringido por
\texttt{math::min(samples*sizeof(FIFO::DATA)+4,\,FIFO::SIZE)}. O valor \artigo{$9088$}
$\geq 1.024$ dispara a reinicialização da primeira guarda, de modo que nenhum
acesso fora dos limites é alcançável.}

\revisao{O candidato surgiu porque o harness omitia essas guardas, e foi
capturado pela etapa obrigatória de \artigo{conferência contra o código de produção}. Este caso é o exemplo mais
didático da tese central deste trabalho:} \artigo{o candidato chegou a ser
proposto \textit{upstream}, como uma limitação incluída no PR~\#27575, antes da
conferência; foi a leitura do driver de produção que mostrou tratar-se de um
contraexemplo formalmente válido para um defeito inexistente.}

\artigo{A refutação foi feita por leitura do código-fonte do driver, e não por
\textit{replay}: para repeti-la, basta conferir as três guardas acima em
\texttt{BMI088\_Accelerometer.cpp}. O artefato depositado não contém programas
de \textit{replay} nem registros de sanitizadores.}

\lc{Refutação depende de leitura; verificar formalmente as três guardas, como feito para o GPS.}

\section{Estudo de Caso: Driver GPS uBlox}
\label{sec:caso-gps}

O \artigo{harness do} driver uBlox (\artigo{arquivo-fonte de $1.712$}~LoC) foi verificado em $64,9$\,s. A verificação de
espaço disponível \artigo{subtrai o nível de preenchimento da capacidade}
(Listagem~\ref{lst:gps-vulneravel}):

\begin{lstlisting}[
  caption={\artigo{Verificação de espaço disponível no \textit{harness} do GPS uBlox}.},
  label={lst:gps-vulneravel},
  numbers=left,
  frame=single,
  language=C]
if (write_len > GPS_DUMP_DATA_SIZE - dump_data->len)
\end{lstlisting}

\artigo{No harness, \texttt{GPS\_DUMP\_DATA\_SIZE} é a constante inteira $200$
(\texttt{\#define}), e o comprimento \texttt{uint8\_t} é promovido a
\texttt{int}; quando \texttt{dump\_data->len} excede a capacidade, a diferença
é negativa ($-7$ para \texttt{len = 207}) e, na comparação com
\texttt{write\_len} (\texttt{size\_t}), converte-se em um valor sem sinal
enorme, tornando a guarda sempre falsa. No driver de produção a expressão
correspondente é \texttt{sizeof(dump\_data->data) - dump\_data->len}
(\texttt{gps.cpp}, linha 741), com buffer de 79 bytes
(\texttt{msg/GpsDump.msg}); ali a subtração já é feita em \texttt{size\_t}, sem
sinal, e daria a volta módulo $2^{32}$ no alvo de voo STM32 de 32 bits ---
\textit{se} \texttt{len} pudesse exceder 79, o que a
Seção~\ref{subsec:gps-inducao} mostra não ocorrer.}

O ESBMC produziu \artigo{\texttt{input\_len = 10} e \texttt{dump\_data->len = 207}
(\texttt{test\_choice = 3}, função \texttt{test\_gps\_real\_bit\_operation};
idêntico nas três execuções)}; o
\texttt{memcpy} subsequente escreve fora do buffer. Classificação:
%% [BANCA - Prof. Yuri] "Isso ai nao e underflow, isso e overflow."
%% DECISAO CONFIRMADA PELO LC: mantida a terminologia da CWE, com nota de
%% rodape reconhecendo a distincao. O CWE-191 do MITRE denomina-se literalmente
%% "Integer Underflow (Wrap or Wraparound)"; trocar o termo desalinharia a
%% dissertacao da classificacao adotada e do artigo submetido.
%% (Caixa \pergunta removida: questao respondida.)
\corr{no \textit{harness}, em que a subtração é feita com sinal e o resultado
negativo é depois convertido para um tipo sem sinal,} \textbf{CWE-195}
\corr{(\textit{Signed to Unsigned Conversion Error})} $\rightarrow$
\textbf{CWE-119} (\artigo{\textit{Improper Restriction of Operations within the
Bounds of a Memory Buffer}}); \corr{na expressão do driver de produção, em que os
dois operandos já são sem sinal, a classe correspondente seria} \textbf{CWE-191}
(\textit{Integer Underflow \artigo{(Wrap or Wraparound)}}).\footnote{\revisao{Registra-se
uma ambiguidade terminológica. No vocabulário de análise numérica,
\textit{underflow} designa o resultado pequeno demais para ser representado, ao
passo que o fenômeno aqui descrito --- uma subtração sem sinal que dá a volta
para um valor muito grande --- corresponderia a \textit{overflow} por
\textit{wraparound}. Adota-se, ainda assim, a denominação da \textit{Common
Weakness Enumeration}, cuja entrada CWE-191 intitula-se \textit{Integer Underflow
(Wrap or Wraparound)} e cobre exatamente esse caso, de modo a preservar a
correspondência entre o texto, a classificação empregada e os relatórios
submetidos à comunidade.}} \corr{A Figura~\ref{fig:ce-gps} reproduz a saída do
verificador para esse contraexemplo.} \corr{O \texttt{dump\_buffer} local do \textit{harness}
corresponde ao buffer \texttt{dump\_data} passado ao driver, e o traço literal
consta do Apêndice~\ref{apend:contraexemplos}.}

\begin{figure}[H]
    \centering
    \definecolor{ceFundo}{HTML}{F7F8FA}\definecolor{ceCinza}{HTML}{5A6070}%
    \definecolor{ceAmarelo}{HTML}{0B5CAD}\definecolor{ceVermelho}{HTML}{B3261E}%
    \begin{tcolorbox}[colback=ceFundo, colframe=black!45, boxrule=0.6pt, arc=2pt,
      left=8pt, right=8pt, top=6pt, bottom=6pt, width=0.95\linewidth,
      fontupper=\ttfamily\scriptsize\linespread{1.45}\selectfont]
    \textcolor{ceVermelho}{\textbf{[Counterexample]}}\\
    \textcolor{ceCinza}{State 1 -{}- test\_gps\_driver.cpp:224 -{}- main}\\
    \hspace*{1em}\fcolorbox{ceAmarelo}{blue!6}{\textcolor{ceAmarelo}{\textbf{test\_choice = 3 (...00000000011)}}}\\
    \textcolor{ceCinza}{State 3 -{}- test\_gps\_driver.cpp:179 -{}- test\_gps\_real\_bit\_operation}\\
    \hspace*{1em}\fcolorbox{ceAmarelo}{blue!6}{\textcolor{ceAmarelo}{\textbf{input\_len = 10 (...00000001010)}}}\\
    \textcolor{ceCinza}{State 8 -{}- test\_gps\_driver.cpp:185 -{}- test\_gps\_real\_bit\_operation}\\
    \hspace*{1em}\fcolorbox{ceAmarelo}{blue!6}{\textcolor{ceAmarelo}{\textbf{dump\_buffer.len = \{ .data=\{...\}, .len=207, ... \}}}}\\
    \textcolor{ceVermelho}{\textbf{Violated property:}}\\
    \textcolor{ceVermelho}{\hspace*{1em}file string.c line 280 function \_\_memcpy\_impl}\\
    \textcolor{ceVermelho}{\hspace*{1em}dereference failure: Access to object out of bounds}\\
    \textcolor{ceVermelho}{\textbf{VERIFICATION FAILED}}\\[-2pt]
    \textcolor{ceCinza}{\rule{\linewidth}{0.3pt}}\\
    \textcolor{ceAmarelo}{harness: 200 - 207 = -7 (int) -> size\_t enorme na compara\c{c}\~ao; guarda sempre falsa}\\
    \textcolor{ceCinza}{ESBMC 8.1.0 + Boolector 3.2.4 -{}- log: medianas\_logs/GPS\_run1.log (igual em run2 e run3)}\\
    \textcolor{ceCinza}{PX4 main, commit 8be2590c15 -{}- Issue \#26866 / PR \#27576}
    \end{tcolorbox}
    \caption[Contraexemplo do GPS uBlox (saída do ESBMC)]{\corr{Saída do ESBMC para o contraexemplo do GPS uBlox, montada a partir do log (ESBMC~v8.1.0, Boolector~3.2.4).}}
    \label{fig:ce-gps}
\end{figure}

Em hardware Pixhawk sem MMU, uma escrita fora dos limites pode corromper
memória adjacente; sob o modelo de memória plana do NuttX ela não
necessariamente levanta uma falha limpa. \revisao{O impacto preciso,
porém, depende do arranjo em tempo de ligação}\artigo{, que não foi analisado
aqui, e da alcançabilidade no alvo, tratada a seguir}\revisao{.}

\artigo{A análise do código-fonte do driver de produção resolve essa questão
negativamente: o laço de escrita limita \texttt{write\_len} ao espaço livre
restante, de modo que o nível de preenchimento nunca excede o buffer antes de o
registro ser publicado e o campo zerado. A violação é, portanto, alcançável no
modelo, mas não no driver de produção. Ela é mantida aqui como a ilustração
mais nítida da tese central da metodologia: um contraexemplo formalmente
válido exige triagem contra o código de produção antes de poder ser chamado de
defeito.} \artigo{Essa conclusão, obtida pelo mantenedor e por leitura do
código, foi depois verificada formalmente (Seção~\ref{subsec:gps-inducao}).}

\begin{tcolorbox}[colback=blue!5!white, colframe=blue!60!black,
title=\textbf{Resposta à QE2 --- Violações em código sem teste unitário direto}]
\corr{Parcial. O BMC expôs, no \textit{harness} do GPS, uma violação de
segurança de memória alcançável no modelo, em uma função
(\texttt{dumpGpsData}) que nenhum teste unitário do PX4 referencia. O estado
violador, porém, é inalcançável no driver de produção
(Seção~\ref{subsec:gps-inducao}); nenhum teste, unitário ou de integração,
poderia exercitá-lo. O que a campanha demonstra para a QE2 é, portanto, mais
restrito: a função não tem cobertura unitária direta, e a análise formal foi o
que estabeleceu tanto a violação no modelo quanto a sua inalcançabilidade na
produção. Na primeira campanha, nenhum defeito do driver foi exposto.}
\end{tcolorbox}

\corr{A ausência de teste unitário direto foi conferida em 09/07/2026 por busca
textual na árvore do PX4 (registro em \texttt{evidencia\_unit\_test\_gps.txt},
no repositório de artefatos). As únicas ocorrências de \texttt{dumpGpsData}
estão em \texttt{src/drivers/gps/gps.cpp} (declaração na linha~303, chamadas nas
linhas~428, 435, 457 e~660, definição na linha~727); nenhum arquivo de teste
inclui \texttt{gps.cpp} nem instancia a classe \texttt{GPS}; e o único teste do
diretório, \texttt{devices/gps-parser-test.cpp}, não menciona a função. O
registro não guarda os comandos exatos nem o \textit{commit}; os números de
linha coincidem com os do \textit{commit} \texttt{8be2590c15}. A busca evidencia
ausência de cobertura unitária direta, e não de exercício indireto por SITL ou
por testes de integração.}

\subsection{Verificação do \textit{patch}}
\label{subsec:patch-gps}

\artigo{O artigo registra uma limitação da verificação do \textit{patch} então
disponível: o CBMC~5.95.1 reportava $0$ de $7$ propriedades violadas no
harness com a guarda do PR~\#27576, mas o harness \textbf{sem} a guarda
reportava o mesmo $0$ de $7$, porque ambos assumiam
\texttt{dump\_buffer.len}~$<$~\texttt{GPS\_DUMP\_DATA\_SIZE} --- pré-condição
que exclui justamente os estados que a guarda trata. A comparação, portanto,
não isolava o efeito da guarda, e não havia log do ESBMC para o harness
corrigido.}

\artigo{Para fechar essa lacuna, a verificação foi refeita em 11/10/2026 com
\textbf{as mesmas pré-condições com e sem o \textit{patch}}, em duas formas da
guarda: (a) a do PR~\#27576 (subtração em \texttt{int}; se não há espaço,
zera \texttt{len} e interrompe o laço); e (b) a forma pedida pelo mantenedor na
revisão do PR (no topo do laço, se \texttt{len} excede a capacidade, zera
\texttt{len} e prossegue). A Tabela~\ref{tab:patch-gps} resume os
desfechos.}

\begin{table}[H]
\centering
\caption{\artigo{Verificação do \textit{patch} do GPS sob as mesmas pré-condições}}
\label{tab:patch-gps}
\footnotesize
\begin{tabular}{|p{6.2cm}|c|c|c|}
\hline
\textbf{Ferramenta e harness} & \textbf{Sem \textit{patch}} & \textbf{PR \#27576} & \textbf{Forma do mantenedor} \\
\hline
\artigo{ESBMC~8.1.0 + Boolector~3.2.4; harness da Tabela~\ref{tab:resultados}, mesmas \textit{flags}} &
\artigo{\textit{violação}} & \artigo{sem violação} & \artigo{sem violação} \\
\artigo{CBMC~5.95.1; janela $\texttt{len} \geq 190$; \texttt{--unwind} 2, 3 e 5} &
\artigo{\textit{violação}} & \artigo{sem violação} & \artigo{sem violação} \\
\hline
\end{tabular}

\footnotesize
\begin{flushleft}
\artigo{ESBMC: sem \textit{patch}, 4.463 VCCs e o mesmo contraexemplo da
Tabela~\ref{tab:resultados} (\texttt{input\_len = 10},
\texttt{dump\_buffer.len = 207}); PR~\#27576, 3.453 VCCs; forma do mantenedor,
4.760 VCCs. CBMC: sem \textit{patch}, 1 de 73 propriedades violadas
(\textit{memcpy destination region writeable}) nos três limites; com as
guardas, 0 de 73 e 0 de 79. As cópias do harness com \textit{patch} diferem do
original apenas no laço de \texttt{dumpGpsData()}. O ESBMC~8.1.0 usado foi
recompilado da \textit{tag} \texttt{v8.1} com Boolector~3.2.4 e validado contra
os logs de julho (mesmas contagens de VCC e mesmos contraexemplos). Execuções
em máquina de nuvem: só os vereditos são reportados; os tempos não são
comparáveis aos da Tabela~\ref{tab:resultados}.}
\end{flushleft}
\end{table}

\artigo{Sob as mesmas pré-condições, as duas formas da guarda eliminam a
violação que o harness sem \textit{patch} exibe, nas duas ferramentas. O
escopo continua limitado: são verificações até o limite de desenrolamento, sem
asserções de desenrolamento, como na Tabela~\ref{tab:resultados}; e, como a
próxima seção mostra, o estado que a guarda trata não ocorre no driver de
produção, de modo que o \textit{patch} só se justifica como endurecimento
defensivo --- exatamente o enquadramento dado pelo mantenedor.}

\subsection{\artigo{Inalcançabilidade no driver de produção: prova indutiva}}
\label{subsec:gps-inducao}

\artigo{O argumento do mantenedor --- que \texttt{len} nunca excede a
capacidade no programa real --- foi verificado formalmente sobre uma cópia
literal do laço de \texttt{GPS::dumpGpsData()} (\texttt{gps.cpp}, linhas
727--761 do \textit{commit} \texttt{8be2590c15}), com os tipos reais:
\texttt{len} em \texttt{size\_t}, \texttt{sizeof(dump\_data->data)}~$= 79$ e
\texttt{dump\_data->len} em \texttt{uint8\_t}. A prova é por indução sobre o
invariante $I$: \texttt{dump\_data->len}~$< 79$.}

\begin{itemize}
\item \artigo{\textbf{Base.} \texttt{GPS::initializeCommunicationDump()} zera os
dois buffers com \texttt{memset} (linhas 717--718), logo $I$ vale antes da
primeira chamada.}
\item \artigo{\textbf{Passo.} A partir de um estado arbitrário que satisfaz $I$
e de um comprimento de entrada arbitrário ($1 \leq \texttt{len} \leq 4096$), uma
iteração do laço não acessa fora dos limites no \texttt{memcpy} e preserva $I$.
Verificado com o CBMC~5.95.1 (0 de 124 propriedades violadas) e com o
ESBMC~8.1.0 + Boolector~3.2.4 (447 VCCs, \textsc{verification successful}, com
asserções de desenrolamento \textbf{ligadas}).}
\item \artigo{\textbf{Premissa.} \texttt{dumpGpsData()} é a única função que
escreve \texttt{dump\_data->len}: no \textit{commit} verificado, as escritas
estão nas linhas 747--757, e os ponteiros para os buffers são membros privados
da classe.}
\end{itemize}

\artigo{Como controle negativo, o mesmo corpo sem a hipótese $I$ é refutado
pelas duas ferramentas (acesso fora dos limites no \texttt{memcpy}), o que mostra
que o harness detectaria o defeito se o invariante não valesse. Diferentemente
dos resultados da Tabela~\ref{tab:resultados}, este não depende de limite de
desenrolamento: o harness não tem laço próprio, e o único laço (o modelo de
\texttt{memcpy}, de no máximo 79 iterações) é coberto com asserção de
desenrolamento ativa. Logo, em toda execução do driver real o \texttt{memcpy}
permanece dentro do buffer, e o estado violador do harness
(\texttt{len = 207}) é inalcançável na produção. O resultado transforma a
triagem por leitura de código em evidência formal e ilustra a etapa~(v) da
metodologia aplicada ao próprio contraexemplo.}

\corr{A premissa de escritor único foi conferida por busca dos dois ponteiros
em \texttt{gps.cpp} no \textit{commit} verificado. \texttt{\_dump\_to\_device} e
\texttt{\_dump\_from\_device} aparecem apenas na declaração (linhas 224--225),
na liberação (410--411), na alocação e no teste de nulo (709--712), no
\texttt{memset} (717--718) e na seleção do buffer dentro de
\texttt{dumpGpsData()} (729). Não são passados a outra função nem expostos por
método de acesso, e a publicação recebe o conteúdo por referência para cópia;
não há, portanto, outro nome pelo qual \texttt{len} possa ser escrito
(\textit{aliasing}).}

\corr{A conclusão depende de três hipóteses, que a prova não verifica: (i)~as
chamadas a \texttt{dumpGpsData()} sobre um mesmo buffer não se intercalam, isto
é, não há escrita concorrente; (ii)~a memória do buffer não é corrompida por um
defeito em outra parte do \textit{firmware}; e (iii)~compilador e
\textit{hardware} executam o código conforme a semântica da linguagem. Sob
qualquer violação dessas hipóteses o invariante pode deixar de valer, e a guarda
proposta no PR~\#27576 passaria a ter efeito.}

\section{Discussão}
\label{sec:discussao-resultados}

\subsection{Configuração, \textit{soundness} e tempo}

\begin{tcolorbox}[colback=blue!5!white, colframe=blue!60!black,
title=\textbf{Resposta à QE3 --- Configurações por módulo}]
\corr{Nove dos doze módulos verificam-se sob a invocação padrão, sem qualquer
ajuste. Apenas três exigiram configuração específica: MAVLink com
\texttt{--unwind~300} (laços profundos de \textit{parsing}), Navigator com
\texttt{--unwind~15} e VTOL com \texttt{--no-library} (desabilita a biblioteca C
abstrata). Todas estão documentadas no repositório de artefatos. O custo de
configuração é, portanto, baixo: a maioria dos módulos não requer
sintonia alguma.}
\end{tcolorbox}

%% LC (resolvido 12/11/2026): o esforço humano não foi registrado; a afirmação
%% de baixo custo foi restrita ao tempo de máquina e à configuração.
\corr{Essa resposta refere-se apenas à configuração do verificador. O esforço de
escrever cada \textit{harness}, triar os contraexemplos e conferir o código de
produção não foi registrado durante a pesquisa, e por isso não é reportado em
horas por módulo; é esse esforço, e não o tempo de máquina, que domina o custo
de adoção (Seção~\ref{sec:conclusao-limitacoes}).}

%% [BANCA - Prof. Barreto] "Ha uma tendencia do pessoal de metodos formais em achar que
%% o teste e pior que a verificacao formal. Cada um tem o seu lugar."
\revisao{Convém evitar uma assimetria comum na literatura de métodos formais: apontar a
cobertura limitada dos testes sem reconhecer que a verificação formal está sujeita à mesma
limitação, apenas por outra via. A cobertura do teste é limitada pelos casos que alguém
escreveu; a do BMC é limitada por três fatores igualmente humanos --- as \emph{propriedades}
que se escolheu especificar, a \emph{profundidade} até a qual se decidiu desenrolar, e o
\emph{recorte} de código submetido ao verificador. Uma falha de lógica de negócio para a qual
nenhuma propriedade foi formulada passa despercebida por um verificador tão silenciosamente
quanto por uma suíte de testes que não a exercita.}

\revisao{A diferença entre as duas abordagens não é, portanto, de completude, mas de
\emph{natureza da garantia}: dentro do recorte especificado, o BMC cobre exaustivamente todas
as entradas possíveis até o limite, enquanto o teste cobre exatamente as entradas escolhidas.
Cada uma alcança o que a outra não alcança, e é nesse sentido preciso que são complementares
--- e não em uma suposta superioridade de uma sobre a outra.}

O BMC é \textit{sound} \artigo{e completo até o limite $k$}: violações reportadas são
reais para o modelo codificado, mas a ausência de violações até a profundidade
$k$ não é prova~\cite{Biere1999}.
Dentro desse escopo, BMC e teste unitário são
complementares~\cite{Godefroid2005}.

\corr{Os dois módulos mais caros da campanha, à parte o GPS, são o MAVLink
($49{,}2$\,s) e o EKF2 ($12{,}2$\,s). No caso do MAVLink, o custo acompanha o
limite de desenrolamento elevado: com $k = 300$, cada iteração replica as
instruções do laço na fórmula SMT, conforme discutido na
Seção~\ref{subsec:escolha-k}. Não se atribui aqui o custo ao \textit{bit-blasting}
de IEEE-754~\cite{Gadelha2020}, porque os dados da Tabela~\ref{tab:resultados}
não sustentam essa explicação: o Expo RC é aritmética de ponto flutuante pura e
verifica-se em $4{,}9$\,s, o Barômetro em $0{,}4$\,s e o VTOL, que emprega
\texttt{sinf} e \texttt{sqrtf}, em $4{,}0$\,s. Isolar a contribuição do ponto
flutuante exigiria execução controlada adicional, registrada como trabalho
futuro.}

\subsection{\revisao{Generalidade do padrão encontrado}}

\revisao{Dois outros candidatos com subtração próxima a \texttt{memcpy}
foram inspecionados (o \textit{statustext} do MAVLink e o RC Update) e não
apresentam \textit{underflow} desprotegido. O padrão do GPS, portanto, não é
disseminado nos módulos inspecionados.}

\lc{Método de inspeção não informado; dizer se ESBMC ou leitura, e listar locais no código.}

\subsection{\revisao{Observação entre autopilotos e integração contínua}}
\label{subsec:ardupilot-ci}

\revisao{A mesma observação de ausência de validação de faixa \artigo{(CWE-1284,
\textit{Improper Validation of Specified Quantity in Input}, subclasse de
CWE-20)} aparece no driver equivalente do ArduPilot~\cite{ArduPilot}: a execução do
CBMC~5.95.1 sobre \texttt{AP\_InertialSensor\_BMI088.cpp} (linhas $383$--$385$)
a reproduz, registrada na Issue~\#33272\footnote{\url{https://github.com/ArduPilot/ardupilot/issues/33272}}.
Como os mantenedores do PX4 recusaram essa mesma observação, ela é reportada
apenas como verificação de consistência entre autopilotos, e \textbf{não} como
achado independente.} \artigo{Um mantenedor do ArduPilot encerrou a
\textit{issue} citando a faixa de medição de $-104$ a $150\,^{\circ}$C do
\textit{datasheet}; os dois projetos, portanto, recusaram a observação.}
\revisao{Vale notar que o ArduPilot já protege o contador do FIFO
(\texttt{fifo\_length}~$\leq 8 \times 7 = 56$ bytes,} \artigo{linhas 314--315 do
ramo \textit{master} consultado em 12/10/2026}\revisao{).}

\revisao{Um fluxo de trabalho do GitHub Actions em um \textit{fork} do
PX4\footnote{\url{https://github.com/marcela2881/PX4-Autopilot}} executa os
harnesses do IMU e do GPS a cada \textit{push} e \textit{pull request},
integrando verificação baseada em harness a um \textit{pipeline} de CI padrão.
Os tempos de relógio da CI não foram medidos sob configuração controlada.}

%% [BANCA - Prof. Yuri] "A metodologia promete avaliar impacto e consequencias da falha
%% no contexto de voo, mas eu nao senti que isso esta sendo feito nos exemplos."
%% [BANCA - Prof. Reginaldo] "Nao e trivial fazer essa analise de criticidade. Para essa
%% temperatura, sendo simples, da ate para comentar."
%% DECISAO CONFIRMADA PELO LC: a analise e feita onde e tratavel (temperatura),
%% com o criterio do Prof. Reginaldo, e declara-se por que nao se estende aos
%% demais. O LC respondeu "vale tentar algo mais ambicioso" quanto ao SITL, que
%% fica registrado como trabalho futuro por nao caber no cronograma ate a defesa.
%% (Caixa \pergunta removida: questao respondida.)
\section{\revisao{Análise de Severidade e Plausibilidade Operacional}}
\label{sec:severidade}

\revisao{Um contraexemplo demonstra que existe uma execução que viola a propriedade,
mas nada diz sobre duas questões que interessam a quem opera o veículo: \emph{quão grave} é a
violação e \emph{quão plausível} é que o estado que a produz ocorra em voo. O verificador não
ordena achados por severidade nem avalia verossimilhança operacional --- essa etapa é humana e
constitui parte da triagem. Esta seção a executa onde ela é tratável e explicita por que não
se estende a todos os casos.}

\subsection{\revisao{Temperatura do BMI088}}

\revisao{Para esta propriedade a análise de severidade é decidível a partir do próprio
código, e reduz-se a uma pergunta: existe, a jusante do driver, algum mecanismo que sinalize a
leitura anômala?}

\revisao{Se a pilha de software emite alarme ao ultrapassar os $85\,^\circ$C
especificados, então o fato de a variável admitir $86{,}25\,^\circ$C é apenas margem de
representação, sem consequência operacional: o sistema já sabe que a leitura é inválida. Se,
ao contrário, nenhuma camada sinaliza a excursão, a situação é qualitativamente distinta ---
o veículo pode estar operando fora da faixa certificada do componente sem que isso seja
comunicado ao operador ou à lógica de \textit{failsafe}. A responsabilidade em jogo não é a de
impedir o sobreaquecimento, que é fenômeno físico, mas a de \emph{medir e informar} a
temperatura, que é responsabilidade de software.}

\revisao{No PX4} \artigo{(\textit{commit} \texttt{8be2590c15})}\revisao{, a temperatura do BMI088 não alimenta o caminho inercial e não
existe alarme dedicado a excursões dessa grandeza no nível do driver. A severidade, portanto,
não é a de perda de controle, mas a de \textbf{perda de observabilidade}: uma leitura fora de
faixa é consumida silenciosamente. Essa é precisamente a razão pela qual a observação foi
mantida como nota de projeto após a recusa upstream --- é um achado de qualidade de
instrumentação, não de segurança de voo, e essa distinção só se torna visível quando a análise
de severidade é feita.}

\corr{Os consumidores da leitura foram levantados por busca textual no
\textit{commit} analisado. O driver publica a temperatura depois de testar
apenas se o valor é finito (\texttt{PX4\_ISFINITE}); não há comparação com a
faixa do \textit{datasheet}. A jusante, o campo é lido em quatro lugares:
\texttt{sensors/vehicle\_imu}, que calcula a média e a publica em
\texttt{vehicle\_imu\_status}; os fluxos \texttt{SCALED\_IMU} do MAVLink, que a
enviam à estação de solo; \texttt{drivers/heater}, que a usa como realimentação
do aquecedor da IMU; e \texttt{modules/temperature\_compensation}, que, quando
habilitado por parâmetro (\texttt{TC\_A\_ENABLE}), calcula a partir dela a
correção de \textit{offset} do acelerômetro. Nenhum desses pontos compara o
valor com a faixa de $-40$ a $+85\,^{\circ}$C. Dois deles qualificam a afirmação
de que a temperatura não alimenta o caminho inercial: com a compensação térmica
habilitada ela influencia a correção aplicada à medida, e no controle do
aquecedor ela fecha uma malha. O efeito de uma leitura fora de faixa nesses dois
caminhos não foi analisado nesta dissertação.}

\subsection{\revisao{Plausibilidade operacional da violação do GPS}}

\revisao{Para o GPS a análise inicialmente ficou incompleta. A
violação exige que \texttt{dump\_data->len} exceda 200 no momento da chamada. Como o campo é
acumulado entre chamadas sucessivas, e como o recurso de \textit{dump} é ativado por
parâmetro, estabelecer plausibilidade exigiria modelar a sequência de chamadas a montante ---
isto é, demonstrar que os módulos abstraídos pelos \textit{stubs} podem efetivamente conduzir
o driver a esse estado.}

\artigo{Essa modelagem foi feita posteriormente, não por simulação da sequência de
chamadas a montante, mas por análise direta do laço de escrita do próprio driver: ele
limita \texttt{write\_len} ao espaço livre restante antes de cada incremento de
\texttt{dump\_data->len}, de modo que o campo nunca excede a capacidade do buffer
antes de o registro ser publicado e zerado. A pré-condição que o \textit{harness} assume
--- \texttt{dump\_data->len} já acima de 200 ao entrar na verificação --- não tem
como surgir no driver de produção, porque é o próprio laço de escrita que mantém
esse invariante.}

\revisao{A severidade, caso o estado fosse atingido, seria alta: em hardware sem
MMU uma escrita fora dos limites pode corromper estado adjacente de voo.}
\artigo{A plausibilidade, no entanto, está resolvida: o estado que dispara a
violação é inalcançável no driver de produção. É essa combinação --- consequência
potencialmente severa, mas restrita ao modelo do \textit{harness} --- que melhor
descreve o achado, e que motivou a mudança de enquadramento junto ao mantenedor,
de vulnerabilidade para endurecimento defensivo (PR~\#27576).}

\artigo{A confirmação por simulação \textit{software-in-the-loop}, cogitada
como caminho para estabelecer plausibilidade, deixou de ser necessária para
este achado específico: a análise do código-fonte a resolveu diretamente. SITL
permanece de interesse para casos futuros em que a leitura do código não for
conclusiva por si só, e fica registrada nesse sentido mais geral entre as
direções de trabalho futuro.}

\subsection{\revisao{Por que a análise não se estende a todos os módulos}}

\revisao{Dos doze módulos, dez não produziram violação, de modo que não há severidade a
avaliar. Entre os três contraexemplos, o candidato do FIFO foi refutado \artigo{na conferência contra o código de produção} e
sua severidade é, por construção, nula. Restam os dois casos analisados acima. A análise de
criticidade em sentido pleno --- estabelecer a relação causal entre a violação e o desfecho da
missão --- exigiria modelo da dinâmica do veículo e da propagação de erro entre módulos,
elemento que está fora do escopo desta dissertação e que se registra como direção de trabalho
futuro (Seção~\ref{sec:trabalhos-futuros-hw}).}

\section{\corr{Resultados da Segunda Campanha}}
\label{sec:resultados-segunda-campanha}

\corr{A segunda campanha aplicou o método a dez módulos complementares, com
verificação sobre o código original do PX4, validação cruzada com o CBMC em
todos os \textit{harnesses} e \textit{replay} sob sanitizadores. O protocolo,
as tabelas por frente e os registros estão na
Seção~\ref{sec:segunda-campanha}; esta seção reúne os desfechos
(Tabela~\ref{tab:desfechos-segunda}) e a validação em bancada.}

\begin{table}[!htb]
\centering
\caption{\corr{Desfechos da segunda campanha, por módulo}}
\label{tab:desfechos-segunda}
\footnotesize
\begin{tabular}{|p{5.6cm}|c|p{7.4cm}|}
\hline
\textbf{Desfecho} & \textbf{Módulos} & \textbf{Quais} \\
\hline
\corr{Defeito reproduzido no código real e corrigido \textit{upstream} por terceiros} & \corr{3} & \raggedright\corr{\texttt{externalChecks}, \texttt{ft7\_technologies}, \texttt{IridiumSBD}} \tabularnewline
\hline
\corr{Defeito reproduzido e reportado neste trabalho} & \corr{2} & \corr{RoboClaw (\textit{Issue}~\#29785), batmon (\textit{Issue}~\#29787)} \\
\hline
\corr{Defeito reproduzido, relatado por terceiro e não corrigido} & \corr{1} & \corr{\texttt{sumd}} \\
\hline
\corr{Defeito reproduzido, não reportado (impacto desprezível)} & \corr{1} & \corr{\texttt{MicroStrain}} \\
\hline
\corr{Falso positivo de modelagem, barrado na conferência} & \corr{1} & \corr{\texttt{bst}} \\
\hline
\corr{Sem defeito de memória} & \corr{2} & \corr{\texttt{logger}, \texttt{Common}} \\
\hline
\end{tabular}
\end{table}

\corr{Em sete dos dez módulos o efeito previsto pelo verificador foi reproduzido
executando o código real do PX4. Em três deles, o próprio projeto corrigiu o
mesmo ponto de forma independente, o que corrobora os achados sem que fosse
preciso reportá-los. O único falso positivo, no \texttt{bst}, foi barrado pela
conferência antes de qualquer relato público. O contraste com a primeira
campanha é direto: lá, três contraexemplos não resultaram em defeito confirmado
algum; aqui, nenhum candidato chegou a ser reportado sem ter sido antes
reproduzido no código real. Os dois relatos abertos aguardavam triagem na data
de fechamento deste texto. A campanha rendeu também uma correção fora do PX4: um
defeito no modelo de \texttt{snprintf} do ESBMC, reportado e corrigido pelos
mantenedores da ferramenta.}

\subsection{\corr{Validação em bancada}}
\label{subsec:validacao-bancada}

\corr{Dois ensaios foram feitos numa Pixhawk~6X~RT com \textit{firmware}
compilado do \textit{commit} analisado. Eles respondem a perguntas opostas: um
mostra que um estado apontado pelo verificador \emph{não} ocorre na placa, e o
outro, que um defeito apontado \emph{ocorre}.}

\corr{\textbf{GPS (primeira campanha).} Com o registro da comunicação do GPS
habilitado (\texttt{GPS\_DUMP\_COMM}~$= 1$), foram coletadas 175 mensagens
\texttt{gps\_dump} em dois registros gravados pela placa em bancada, analisadas com
\texttt{pyulog}. Em todas, o campo \texttt{len} vale 207, isto é, o bit~7
(que marca dado enviado ao dispositivo) somado a 79~bytes, a capacidade do
buffer: nenhuma mensagem excede a capacidade. O número coincide por acaso com o
\texttt{len}~$= 207$ do contraexemplo, mas tem outro significado: no
\textit{harness}, 207 é um nível de preenchimento acima de um buffer de
200~bytes; na placa, é um buffer de 79~bytes exatamente cheio. O ensaio é
consistente com a prova da Seção~\ref{subsec:gps-inducao}. Seu limite é que só
o sentido da placa para o receptor foi exercitado, sem receptor GPS conectado.}

\corr{\textbf{RoboClaw (segunda campanha).} Com um ESP32 no papel do
controlador de motores, o valor de \textit{failsafe} em 257 e a terminação de
voo acionada, a placa transmitiu o byte de dados 130, fora da faixa de 0 a 127
aceita pelo protocolo; com o valor em 254, transmitiu 127; e, com a correção
aplicada, voltou a transmitir 127 mesmo com 257. Os pacotes, as contagens e os
limites do ensaio (controlador emulado, \textit{firmware} com o driver
habilitado e dois membros inicializados, apenas o primeiro dos dois defeitos)
estão na Seção~\ref{subsec:2c-roboclaw}.}

\corr{Os dois ensaios cobrem, cada um, um único achado. Os demais resultados
desta dissertação apoiam-se na verificação formal e, na segunda campanha, no
\textit{replay} do código real em computador.}

\section{Ameaças à Validade}
\label{sec:ameacas-validade}

\subsection{Validade interna}

Os tempos reportados são medianas de três execuções e devem ser lidos apenas
como ordens de grandeza. \corr{Na fase exploratória, um resultado do MAVLink com \texttt{--unwind~10}
foi descartado como artefato de configuração: era desenrolamento insuficiente,
resolvido com \texttt{--unwind~300} e confirmado pela reexecução com as
asserções ligadas (Seção~\ref{subsec:escolha-k}). Um segundo candidato inicial,
no EKF2, foi descartado na mesma fase, mas não há registro que permita
descrevê-lo; por isso ele não é contado nem discutido neste trabalho.}
\revisao{O \textit{underflow} do GPS é descarregado apenas pelo
Boolector dentro do \textit{timeout}, de modo que o achado é sensível ao
solver; a validação cruzada com CBMC mitiga artefatos específicos de
ferramenta.}
\artigo{Com o Z3, o ESBMC não produziu veredito dentro do \textit{timeout} no
GPS, nem no MAVLink e no EKF2 com \texttt{--unwind 10}, nas re-execuções com a
versão 8.1 (\texttt{validacao\_esbmc\_8.1/} no repositório); e a validação
cruzada com o CBMC é parcial no GPS, pois só termina com limites reduzidos.}

\corr{O limite de desenrolamento também é ameaça à validade interna.} A verificação foi limitada a $k = 10$ e, com
\texttt{--no-unwinding-assertions} habilitado, laços que excedem o limite são
truncados silenciosamente --- de modo que um veredito \textsc{successful} pode
ocultar falsos negativos mesmo abaixo da profundidade operacional pretendida.
\artigo{A reexecução com as asserções ligadas mostrou que isso não ocorre nos dez
módulos sem violação (Tabela~\ref{tab:unwinding}); no GPS, porém, o laço do
\texttt{memcpy} excede $k = 10$, e a análise desse módulo não o cobre por
inteiro. Permanece, em todos os casos, o limite do recorte: o veredito vale
para o \textit{harness}, e não para o módulo inteiro.}
Os totais de VCC por módulo misturam limites de desenrolamento (10, 15 e 300) e
não são diretamente comparáveis.

\subsection{Validade externa}

Selecionar módulos por viabilidade de \textit{stubbing} sub-representa módulos
com acoplamento pesado a hardware ou ao escalonador, de modo que a amostra não
é estatisticamente representativa do PX4; a mitigação é cobrir as principais
classes de entrada não confiável. A análise modular não captura bugs de
integração decorrentes de interações concorrentes, e os barramentos CAN/UAVCAN,
a ponte RTPS/uXRCE-DDS e os drivers de fluxo óptico e telêmetro permanecem não
verificados.

\subsection{Validade de construto}

\corr{A questão de construto é se as propriedades verificadas medem o que se
pretende afirmar. Três ressalvas se aplicam. Primeiro, as asserções foram
escritas pela autora a partir de \textit{datasheets} e da documentação, e uma
delas, a faixa de temperatura do BMI088, foi contestada pelos mantenedores como
requisito de software: violar a asserção não equivale a violar um requisito do
projeto. Segundo, a asserção final do \textit{harness} do GPS é tautológica, e a
propriedade efetivamente verificada ali é a checagem automática de limites do
\texttt{memcpy} (Capítulo~\ref{cap:metodologia}). Terceiro, toda propriedade é
verificada sobre o \textit{harness}, que transcreve e adapta o código: o
veredito mede a segurança do modelo, e só a conferência contra o código de
produção permite estendê-lo ao driver.}
\corr{A comparação entre ESBMC com Boolector e CBMC com MiniSAT varia
ferramenta e \textit{backend} conjuntamente, de modo que a propriedade distinta
que cada uma reporta sobre o defeito do GPS reflete diferença de
\textit{toolchain}, e não de SMT versus SAT isoladamente.}

\subsection{\corr{Validade de conclusão}}

\corr{Cada módulo foi executado três vezes e reporta-se a mediana, sem medida de
dispersão. Isso sustenta afirmações de ordem de grandeza --- ``menos de noventa
segundos'' --- mas não comparações finas entre módulos cujos tempos diferem por
poucos segundos. Não há, além disso, \textit{baseline} de comparação: nenhuma
outra técnica foi executada sobre os mesmos doze módulos para contraste, à
exceção do CBMC em dois casos. Segue-se que a afirmação de viabilidade é
absoluta, e não relativa: demonstra-se que a verificação é barata em termos
próprios, não que seja mais ou menos barata que alternativas.}

\subsection{\corr{Construção dos \textit{harnesses} com apoio de ferramenta de IA}}

\corr{Conforme registrado na Declaração de Uso de IA, os \textit{harnesses}
foram inicialmente montados com apoio de assistente de IA, a partir das regras
de instrumentação do ESBMC e do código-fonte de cada módulo, e em seguida
revisados manualmente pela autora. Trata-se de ameaça legítima à validade e
merece registro explícito: um \textit{harness} gerado dessa forma pode
reproduzir padrões plausíveis porém incorretos --- em particular, \textbf{omitir
guardas presentes no código de produção}, que foi exatamente o mecanismo do
falso positivo do contador de FIFO discutido na Seção~\ref{subsec:imu-fifo}.}

\corr{A mitigação praticada tem duas camadas. A primeira é a revisão manual de
cada \textit{harness} contra o código-fonte do módulo antes da execução. A
segunda, e mais importante, é a etapa obrigatória de \artigo{conferência} contra o
driver não instrumentado \artigo{(por \textit{replay} ou leitura do código)}, que capturou o único caso em que essa ameaça se
materializou. O episódio do FIFO é, nesse sentido, evidência de que o método
dispõe de mecanismo de refutação embutido, e não apenas de detecção.}

\subsection{\revisao{Correção dos \textit{stubs}}}

\revisao{Os \textit{stubs} que retornam \texttt{nondet\_T()}
sobreaproximam apenas as \textbf{entradas} de hardware, não a lógica do driver
ao redor: nenhuma violação é perdida por conta da abstração de entrada, mas
quando o harness omite guardas de produção --- como mostrou o candidato do FIFO
--- o modelo admite estados que o driver real rejeita. A legitimidade de um
achado repousa, portanto, sobre a etapa obrigatória de \artigo{conferência} contra o
driver não instrumentado.}

\revisao{Inversamente, os \textit{stubs} \textbf{sub}aproximam o
comportamento interno do hardware (máquinas de estado I2C e SPI, escalonamento
do NuttX); defeitos disparados por esse estado interno não seriam detectados, o
que torna a validação com hardware no laço um complemento necessário.}

\subsection{\corr{Defeitos no verificador ou no solver}}

\corr{Um veredito depende da correção do próprio verificador, do seu modelo das
bibliotecas e do solver. A ameaça não é hipotética: durante a segunda campanha,
o modelo de \texttt{snprintf} do ESBMC produziu um alarme falso, depois
confirmado e corrigido pelos mantenedores da ferramenta
(Seção~\ref{sec:segunda-campanha}). A mitigação praticada é parcial. Na
primeira campanha, o CBMC, com outro \textit{front-end} e outro solver, foi
executado apenas sobre os dois módulos com contraexemplo; os dez módulos sem
violação foram verificados com uma única ferramenta e um único solver, de modo
que um defeito que suprimisse uma violação neles passaria despercebido. Na
segunda campanha, a validação cruzada com o CBMC cobre todos os
\textit{harnesses}, e os candidatos alcançáveis foram ainda reproduzidos fora
dos verificadores, por \textit{replay} sob sanitizadores.}

\section{\corr{Lições Aprendidas}}
\label{sec:licoes}

\corr{Quatro lições atravessam as duas campanhas.}

\corr{\textbf{Fidelidade do \textit{harness}.} Os dois falsos positivos deste
trabalho, o contador de FIFO do BMI088 e o \texttt{bst}, têm a mesma origem: o
\textit{harness} omitia guardas que o código de produção contém. Transcrever
trechos é rápido, mas o veredito vale para a cópia. A segunda campanha passou a
verificar o código original e a exigir \textit{replay} da função real, e foi
essa exigência que barrou o \texttt{bst}.}

\corr{\textbf{Reporte prematuro.} Na primeira campanha, as \textit{issues} do
GPS e do FIFO foram abertas antes da conferência contra o código de produção, e
a conferência depois enfraqueceu ou refutou os dois candidatos. A regra adotada
em seguida foi reportar somente depois da conferência, do \textit{replay}, da
consulta à especificação com página citada, da busca por relatos duplicados e
da verificação do estado atual do código. O caso do RoboClaw acrescentou um
item: conferir também em qual ramo o código analisado ainda existe, pois o
driver havia sido reescrito no ramo principal entre a análise e o relato.}

\corr{\textbf{Projeto das hipóteses de entrada.} A pré-condição decide o que o
verificador pode encontrar, nos dois sentidos. No GPS, admitir um estado inicial
acima da capacidade produziu uma violação que o driver real não alcança. No
RoboClaw, restringir a saída à faixa declarada no \texttt{module.yaml}, de 0 a
257, foi o que expôs o defeito; uma faixa mais estreita, de aparência razoável,
o teria ocultado. Cada pré-condição precisa de origem documentada no código ou
na especificação.}

\corr{\textbf{Limites da ferramenta.} Três comportamentos do verificador
exigiram atenção: o truncamento silencioso de laços com
\texttt{--no-unwinding-assertions}, a sensibilidade do contraexemplo do GPS ao
solver e um defeito no modelo de uma função de biblioteca. Nenhum invalidou os
resultados, mas todos reforçam que um veredito isolado, de uma única ferramenta
e configuração, é ponto de partida para a triagem, e não conclusão.}

